% File example.tex
% Contact: simonnet@ecole.ensicaen.fr
%
% version 1.0 (July 24, 2009)
% version 1.1 (September 3, 2009)
% add using of the optional command: \secondauthoraddr

\documentclass[10pt]{article}

\input{includes/icdp2009template.tex}

%other package

% vectorial font
\usepackage{lmodern}

\usepackage{graphicx}
\usepackage{times}


\begin{document}
\noindent

% This should produce references in the order they appear
\bibliographystyle{ieeetr}

\title{Instructions for Preparing and Transferring Final Papers to \\ICPRS 2018}

\authorname{A.B. Author*, C.D. Author+}
\authoraddr{*Affiliation, Country and contact details}

%optional
%\secondauthoraddr{+Affiliation,Country and contact details }


\maketitle

\keywords
Maximum 5 keywords placed before the abstract.

\abstract
This is where the abstract should be placed. It should consist
of one paragraph and a concise summary of the material
discussed in the article below. It is preferable not to use
footnotes in the abstract or the title. The acknowledgement
for funding organisations etc. is placed in a separate section at
the end of the text. We wish you success with the preparation
of your manuscript.

\section{Introduction}
The Proceedings produced for ICPRS 2018 will contain all the
papers accepted for the conference, of which camera-ready versions
need to have been received
by the organisers no later than 9 April 2018. Papers
received after that date will not be included. Papers, which
appear in the Proceedings, have undergone the review process
of the Technical Programme Committee of ICPRS 2018.
Comments of reviewers are made available to the authors. Authors are
asked to prepare and submit electronic versions of their full
papers according to the instructions below.

\section{Manuscript preparation}
Full papers must be typed in English. This instruction page is
an example of the format and font sizes to be used. MS Word
users can download from the conference site these
instructions in Word format.

These are detailed instructions valid for any word
processor. In the title of the paper the initial letters should be
capitalised in all words except articles and prepositions (e.g.:
in, a, an, and, the, there, their, do, on, of, from, with, at etc.).
E.g. "ErDoped Si Nanocrystals as a Candidate for Optical
Amplification" The type should be boldface 18pt and centred
on the page. The author’s name is typed in capital and lower
case bold letters and centred on the page. Directly under the
author’s name in capital and lower case letters and also
centred is the author’s affiliation, address, plus email
address of (at least) the corresponding author. Manuscripts must be
typed single spaced using 10 point characters. Only Times,
Times Roman, Times New Roman and Symbol fonts are
accepted. The text must fall within a frame of 18 cm x 24 cm
centred on an A4 page (21 cm x 29.7 cm).Paragraphs are
separated by 6 points and with no indentation. The text of the
full papers is written in two columns and justified. Each
column has a width of 8.8 cm and the columns are separated
by a margin of 0.4 cm. The maximum length of the full paper
is 6 pages. Do not number the pages. The final format in
which the papers will appear on the Proceedings will be a
PDF file. Authors are required to send a PDF file of their
final paper to be included directly in the Proceedings. All
PDF files should NOT be locked and all fonts and
graphics should be embedded.

\subsection{Figures and tables}
Figures and tables should be centred in the column, numbered
consecutively throughout the text, and each should have a
caption underneath it (see for example Table 1). Care should
be taken that the lettering is not too small. All figures and
tables should be included in the electronic versions of the full
paper. We cannot guarantee that any printed version of the
proceedings will use colour.


\begin{figure}[h]
\centering
\includegraphics[width=1.5cm]{images/fig1}
\caption{\label{tab1}This is an example of a figure caption.} 
\end{figure}


\begin{table}[h]
\begin{center}

\begin{tabular}{c}
nn!1 \\
2 \\
31 \\
6 \\
\end{tabular}
\end{center}
\caption{\label{tab1}This is an example of a table caption.}
\end{table}

\subsection{Equations}
Equations should be typed within the text, centred, and should
be numbered consecutively throughout the text. They should
be referred to in the text as Equation (n). Their numbers
should be typed in parentheses, flush right, as in the following
example.
\begin{equation}
	    PA + A'P - PBR^{-1}B'P + Q  =  0 \enspace.
\end{equation}

\section{Generating a {PDF} file}
The PDF format will be the final format under which the
papers will appear in the Proceedings. Therefore you are
required to submit your paper as PDF document. If this is not
possible, Postscript format is also accepted as long as no fonts
other than the recommended fonts are used.

You can use any of the popular free LaTeX editors (e.g. Kile).

\section{Electronic submission of the full paper}
The submission process for ICPRS 2018 should be done on
line at http://www.icprs.org

A PDF version of your final paper is required. It should
be expected that after your submission, your paper is
published directly from the file you send without any further
proofreading. Therefore, it is advisable for the authors to
print a hard copy of their final version and read it carefully.

\section{Your References}
The list of references should be ordered in the same order as
first cited in the text. All references should be cited in the
text, and using square brackets such as \cite{ref01} and \cite{ref01,ref02}. We
recommend the use of IEEE Transactions style for references.

\section*{Acknowledgements}
The acknowledgement for funding organisations etc. should
be placed in a separate section at the end of the text.



Thank you for your cooperation in complying with these
instructions.


\bibliography{IEEEabrv,includes/icdp2009}


\end{document}
